% formal/library/geom.tex — geom module
% Mathematical results supporting library/src/geom/
%
% Labels defined here:
%   def:symplectic-form, def:J0, def:ehz-capacity, def:lagrangian-product,
%   def:symplectic-product, def:cross-product-4d, def:face-lattice,
%   def:polar-body, def:polytope-dual, def:polygon-h-rep, def:polygon-area,
%   def:reeb-vector-field, def:volume, def:systolic-ratio,
%   lem:piecewise-linear-reeb, lem:positive-span, lem:bounded-triples,
%   lem:rational-pipeline, lem:vertex-enumeration, lem:irredundancy,
%   prop:integer-cramer, lem:shoelace,
%   thm:hko-counterexample,
%   prop:capacity-symplectic-product,
%   prop:prefilter-bound, cor:prefilter-soundness,
%   fact:prefilter-round, fact:prefilter-cast, fact:prefilter-svd,
%   fact:prefilter-solve, fact:prefilter-weyl, fact:prefilter-banach,
%   fact:prefilter-dot, rem:prefilter-constants
%
% Sources: .rs doc comments, Jörn + agent session notes.
%   Polytope definition: def:polytope in thesis (Jörn-approved).
%   Note: formal library files must NOT reference thesis/ files (convention).

% input by formal/library/main.tex and formal/main.tex — do not add \documentclass or \begin{document}

% ====================================================================
% Symplectic structure
% ====================================================================

\begin{definition}[Standard symplectic structure]\label{def:symplectic-form}
% Jörn: text approved (1aac085) — from \begin{definition} to \end{definition}
Write $\mathbb{R}^4$ with coordinates $(q_1, q_2, p_1, p_2)$.
The \emph{standard complex structure} is
\[
  J_0 = \begin{pmatrix} 0 & -I_2 \\ I_2 & 0 \end{pmatrix},
\]
i.e., $J_0(q, p) = (-p, q)$,
where $I_2$ is the $2 \times 2$ identity matrix.
The \emph{standard symplectic form} is
\[
  \omega_0(u, v) \coloneqq \langle J_0 u, v \rangle
  = \sum_{i=1}^{2} (u_{q_i} v_{p_i} - u_{p_i} v_{q_i}),
\]
equivalently $\omega_0 = dq_1 \wedge dp_1 + dq_2 \wedge dp_2$.
\end{definition}


\begin{unverified}
\begin{definition}[Standard complex structure $J_0$]\label{def:J0}
% [TODO: JÖRN - verify: J0 is defined within def:symplectic-form above.
%   This label exists as a separate cross-reference target for code that
%   references J0 specifically (e.g. symplectic_form.rs). Consider whether
%   this should remain a separate label or be merged into def:symplectic-form.]
The standard complex structure $J_0 \in \mathbb{R}^{4 \times 4}$ is
\[
  J_0 = \begin{pmatrix} 0 & -I_2 \\ I_2 & 0 \end{pmatrix},
\]
satisfying $J_0^2 = -I_4$, $J_0^{-1} = J_0^\top = -J_0$,
$\omega_0(J_0 u, J_0 v) = \omega_0(u,v)$, and $|J_0 u| = |u|$.
In the 2D case, $J_0^{(2)} = \bigl(\begin{smallmatrix} 0 & -1 \\ 1 & 0 \end{smallmatrix}\bigr)$.
\end{definition}
\end{unverified}


% ====================================================================
% Capacity and systolic ratio
% ====================================================================

\begin{definition}[EHZ capacity]\label{def:ehz-capacity}
% Jörn: text approved (1aac085) — from \begin{definition} to \end{definition}
The \emph{Ekeland--Hofer--Zehnder capacity} of a convex body $K$ with smooth
boundary is
\[
  c_{\mathrm{EHZ}}(K) \coloneqq \min\bigl\{ A(\gamma) :
  \gamma \text{ is a closed characteristic on } \partial K \bigr\}.
\]
The minimum exists and is positive.
The definition extends to all convex bodies by approximation;
for polytopes, the minimum is attained by a generalized Reeb orbit.
\end{definition}


\begin{unverified}
\begin{definition}[Systolic ratio]\label{def:systolic-ratio}
% [TODO: JÖRN - verify statement]
% Source: dataset.rs doc comment
The \emph{systolic ratio} of a convex body $K \subset \mathbb{R}^4$ is
\[
  \operatorname{sys}(K) \coloneqq \frac{c_{\mathrm{EHZ}}(K)^2}{2\,\operatorname{vol}(K)}.
\]
Viterbo's conjecture asserts $\operatorname{sys}(K) \leq 1$ for all convex bodies $K$.
\end{definition}
\end{unverified}


% ====================================================================
% Products
% ====================================================================

\begin{unverified}
\begin{definition}[Lagrangian product]\label{def:lagrangian-product}
% Source: lagrangian_product.rs doc comments
Let $K_q \subset L_q$ and $K_p \subset L_p$ be convex polygons
with $0 \in \operatorname{int}(K_q)$ and $0 \in \operatorname{int}(K_p)$,
where $L_q = \{(q_1, q_2, 0, 0)\}$ and $L_p = \{(0, 0, p_1, p_2)\}$
are complementary Lagrangian subspaces.
The \emph{Lagrangian product} is the polytope
\[
  K \coloneqq K_q \times K_p
  = \{ (q, p) \in \mathbb{R}^4 : q \in K_q,\; p \in K_p \}.
\]
\end{definition}
\end{unverified}


\begin{unverified}
\begin{definition}[Symplectic product]\label{def:symplectic-product}
% [TODO: JÖRN - verify statement]
% Source: known_polytopes.rs doc comment
Let $A, B \subset \mathbb{R}^2$ be convex bodies containing
the origin in their interiors, where we view the first copy
of~$\mathbb{R}^2$ as the symplectic plane $(q_1, p_1)$ and
the second as $(q_2, p_2)$.
The \emph{symplectic product} is
\[
  A \times_S B
    \coloneqq \{(q_1, q_2, p_1, p_2) \in \mathbb{R}^4
      : (q_1, p_1) \in A,\; (q_2, p_2) \in B\}.
\]
\end{definition}
\end{unverified}


\begin{unverified}
\begin{proposition}[Capacity of symplectic products]%
\label{prop:capacity-symplectic-product}
% Source: standard result in symplectic geometry (see bibliography)
Let $A, B \subset \mathbb{R}^2$ be convex bodies containing
the origin in their interiors, viewed as in
Definition~\ref{def:symplectic-product}.
Then
\[
  c_{\mathrm{EHZ}}(A \times_S B)
    = \min\!\bigl(\operatorname{area}(A),\;
                   \operatorname{area}(B)\bigr).
\]
\end{proposition}

\begin{proof}
We reduce to a product of disks via an area-preserving
map, then squeeze between a ball and a cylinder.

On each copy of~$\mathbb{R}^2$, the symplectic form
restricts to $dq_i \wedge dp_i$, which is the standard
area form, so area-preserving diffeomorphisms are
symplectomorphisms.

% [GAP - JÖRN: the step c_EHZ(A) = area(A) for convex bodies A in R^2
%   was agent-written. The previous version cited "Dacorogna and Moser (1990)"
%   and claimed a "divergence-free vector field extends by zero" to produce a
%   global area-preserving diffeomorphism of R^2. Jörn does not recognize this
%   construction. The entire argument for c_EHZ(A) = area(A) is UNVERIFIED
%   agent-generated math. Needs a correct proof or literature citation.]
We first show $c_{\mathrm{EHZ}}(A) = \operatorname{area}(A)$
for every convex body~$A$ in~$\mathbb{R}^2$.
This follows from the existence of a global area-preserving
(hence symplectic) diffeomorphism mapping~$A$ to a disk
$B^2(r)$ with $\pi r^2 = \operatorname{area}(A)$,
together with symplectic invariance and normalization.

Now let $r_A, r_B > 0$ satisfy
$\pi r_A^2 = \operatorname{area}(A)$ and
$\pi r_B^2 = \operatorname{area}(B)$.
The above gives global symplectomorphisms
$\varphi_1, \varphi_2 \colon \mathbb{R}^2 \to \mathbb{R}^2$
mapping $A$ to~$B^2(r_A)$ and $B$ to~$B^2(r_B)$
respectively.
Since $\omega_0 = dq_1 \wedge dp_1 + dq_2 \wedge dp_2$
decomposes over the two planes, the product
$\varphi_1 \times \varphi_2$ is a symplectomorphism
of~$\mathbb{R}^4$.
By symplectic invariance,
\[
  c_{\mathrm{EHZ}}(A \times_S B)
    = c_{\mathrm{EHZ}}\bigl(B^2(r_A) \times_S B^2(r_B)\bigr).
\]
Assume without loss of generality that $r_A \leq r_B$.
Then
\[
  B^4(r_A)
    \;\subset\; B^2(r_A) \times_S B^2(r_B)
    \;\subset\; B^2(r_A) \times \mathbb{R}^2,
\]
where the first inclusion holds because $r_A \leq r_B$
and the second because $B^2(r_B) \subset \mathbb{R}^2$.
By monotonicity and normalization,
\begin{align*}
  \pi r_A^2
    &= c_{\mathrm{EHZ}}(B^4(r_A)) \\
    &\leq c_{\mathrm{EHZ}}\bigl(B^2(r_A) \times_S B^2(r_B)\bigr) \\
    &\leq c_{\mathrm{EHZ}}(B^2(r_A) \times \mathbb{R}^2)
    = \pi r_A^2.
\end{align*}
Hence $c_{\mathrm{EHZ}}(A \times_S B)
  = \pi r_A^2
  = \min(\operatorname{area}(A), \operatorname{area}(B))$.
\end{proof}
\end{unverified}


% ====================================================================
% Polytope geometry
% ====================================================================

\begin{unverified}
\begin{definition}[Polytope dual / polar body]\label{def:polytope-dual}
% [TODO: JÖRN - verify statement]
% Source: polytope.rs doc comment
Let $K \subset \mathbb{R}^4$ be a polytope with $0 \in \operatorname{int}(K)$
given in $H$-representation $K = \{x : \langle a_i, x\rangle \leq 1\}_{i=1}^F$.
The \emph{polar body} (or \emph{dual}) is
\[
  K^\circ \coloneqq \{y \in \mathbb{R}^4 : \langle x, y \rangle \leq 1
    \text{ for all } x \in K\}
  = \operatorname{conv}(a_1, \ldots, a_F).
\]
\end{definition}
\end{unverified}

% Agent note: def:polar-body used in some .rs files is an alias
% for def:polytope-dual. The label below provides this as a cross-reference target.
\phantomsection\label{def:polar-body}% alias for cross-references from .rs files


\begin{unverified}
\begin{definition}[Face lattice / skeleton]\label{def:face-lattice}
% [TODO: JÖRN - verify statement]
% Source: skeleton.rs doc comment
The \emph{skeleton} (or \emph{face lattice}) of a 4D convex polytope $K$
is the combinatorial structure recording all faces of $K$
(vertices, edges, 2-faces, facets) and their incidence relations.
Two facets $F_i, F_j$ are \emph{adjacent} if $F_i \cap F_j$ is a
face of dimension $\geq 1$.
\end{definition}
\end{unverified}


\begin{unverified}
\begin{definition}[Cross product in $\mathbb{R}^4$]\label{def:cross-product-4d}
% [TODO: JÖRN - verify statement]
% Source: cross_product_4d.rs doc comment
The \emph{cross product} of three vectors $u, v, w \in \mathbb{R}^4$
is the unique vector $u \times v \times w \in \mathbb{R}^4$
perpendicular to all three, with magnitude equal to the volume
of the parallelepiped spanned by $u, v, w$, and orientation
determined by the right-hand rule.
Equivalently,
\[
  (u \times v \times w)_i
  = (-1)^{i+1} \det\bigl(M_i\bigr),
\]
where $M_i$ is the $3 \times 3$ matrix obtained by deleting
row~$i$ from the $4 \times 3$ matrix $(u \mid v \mid w)$.
\end{definition}
\end{unverified}


% ====================================================================
% Polygon geometry
% ====================================================================

\begin{unverified}
\begin{definition}[Convex polygon $H$-representation]\label{def:polygon-h-rep}
% [TODO: JÖRN - verify statement]
% Source: polygon.rs doc comment
A 2D \emph{convex polygon in $H$-representation} is specified by
outward unit normals $n_1, \ldots, n_m \in S^1$ and positive
heights $h_1, \ldots, h_m > 0$:
\[
  P = \bigcap_{i=1}^{m} \{x \in \mathbb{R}^2 : \langle n_i, x \rangle \leq h_i\}.
\]
\end{definition}
\end{unverified}


\begin{unverified}
\begin{definition}[Polygon area]\label{def:polygon-area}
% [TODO: JÖRN - verify statement. "Shoelace formula" here means the standard
%   2D polygon area formula (sum of cross products), NOT lem:shoelace (which is
%   the symplectic action integral). Clarify or rename to avoid ambiguity.]
% Source: polygon.rs doc comment
The \emph{area} of a convex polygon $P$ given in $H$-representation
(Definition~\ref{def:polygon-h-rep}) is computed via triangulation
from the centroid or via the standard 2D shoelace formula on the vertex list.
\end{definition}
\end{unverified}


% ====================================================================
% Volume
% ====================================================================

\begin{unverified}
\begin{definition}[Volume of a 4D polytope]\label{def:volume}
% [TODO: JÖRN - verify statement]
% Source: volume.rs doc comments
The \emph{volume} of a 4D convex polytope $K$ is the standard
Lebesgue measure $\operatorname{vol}(K) = \int_K dx$.
For a 4-simplex $\Delta = \operatorname{conv}(v_0,v_1,v_2,v_3,v_4)$,
\[
  \operatorname{vol}(\Delta)
  = \frac{1}{4!}\,
    \left|\det(v_1 - v_0 \mid v_2 - v_0 \mid v_3 - v_0 \mid v_4 - v_0)\right|.
\]
\end{definition}
\end{unverified}

\begin{unverified}
\begin{lemma}[Origin-star triangulation of volume]\label{lem:volume-star-triangulation}
% [TODO: JÖRN - verify statement and proof]
% Source: volume.rs doc comments
Let
\[
  K = \{x \in \mathbb{R}^4 : \langle a_i, x \rangle \leq 1 \text{ for } i=1,\ldots,F\}
\]
be a bounded polytope. Then $0 \in \operatorname{int}(K)$. For each facet
$F_i$, let $\mathcal{T}_i$ be any triangulation of $F_i$ into tetrahedra.
Then
\[
  \operatorname{vol}(K)
  = \sum_{i=1}^F \sum_{\tau \in \mathcal{T}_i}
    \operatorname{vol}\bigl(\operatorname{conv}(0,\tau)\bigr).
\]
\noindent
In the implementation, $\mathcal{T}_i$ is built by taking the arithmetic
mean of the vertices of $F_i$ and fanning over the ridge polygons in cyclic
order.
\end{lemma}

\begin{proof}
Because $\langle a_i, 0 \rangle = 0 < 1$ for every~$i$, the origin lies in
the interior of~$K$. The cones $\operatorname{conv}(0,F_i)$ over the facets
cover~$K$, and distinct cones intersect only along lower-dimensional faces,
so their overlaps have measure zero. Each facet triangulation $\mathcal{T}_i$
partitions $F_i$ into tetrahedra with pairwise lower-dimensional overlaps.
Coning every tetrahedron in $\mathcal{T}_i$ to~$0$ therefore partitions
$\operatorname{conv}(0,F_i)$ into 4-simplices up to measure-zero overlaps.
Summing over all facets yields the claimed formula.
\end{proof}
\end{unverified}


% ====================================================================
% Reeb dynamics on polytopes
% ====================================================================
% Parameterization: K = {x : a_i^T x ≤ 1}, dual vertices a_i

\begin{unverified}
\begin{definition}[Reeb vector field on polytopes]\label{def:reeb-vector-field}
% [TODO: JÖRN - verify statement]
% Source: reeb_trajectory.rs doc comment
Let $K = \{x : \langle a_i, x \rangle \leq 1\}$ be a polytope
with dual vertices $a_1, \ldots, a_F$.
On facet $F_i$, the \emph{Reeb vector} is
\[
  R_i = 2\, J_0\, a_i.
\]
The Reeb vector $R_i$ is tangent to $F_i$
($\langle a_i, R_i \rangle = 0$, since
$\langle a_i, J_0 a_i \rangle = \omega_0(a_i, a_i) = 0$).
\end{definition}
\end{unverified}


\begin{unverified}
\begin{lemma}[Piecewise-linear Reeb trajectories]\label{lem:piecewise-linear-reeb}
% [TODO: JÖRN - verify statement, add proof]
% Source: reeb_trajectory.rs doc comment
Reeb trajectories on polytope boundaries are piecewise linear:
on each facet $F_i$, the trajectory moves with constant velocity $R_i$.
At a breakpoint where the trajectory transitions from $F_i$ to $F_j$,
the velocity changes discontinuously from $R_i$ to $R_j$.
\end{lemma}
\end{unverified}


% ====================================================================
% Computational geometry lemmas
% ====================================================================

\begin{unverified}
\begin{lemma}[Positive span implies boundedness]\label{lem:positive-span}
% [TODO: JÖRN - verify statement and proof]
% Source: check_bounded_f64_first in vertex_enumeration/boundedness.rs
Let $a_1, \ldots, a_F \in \mathbb{R}^4 \setminus \{0\}$ and
$K = \{x \in \mathbb{R}^4 : \langle a_i, x \rangle \leq 1 \;\forall i\}$.
Then $K$ is bounded if and only if $\{a_1, \ldots, a_F\}$ positively
spans~$\mathbb{R}^4$, i.e., for every $d \in \mathbb{R}^4 \setminus \{0\}$
there exists $i$ with $\langle a_i, d \rangle > 0$.
\end{lemma}

\begin{proof}
($\Leftarrow$) Suppose $\{a_i\}$ positively spans $\mathbb{R}^4$ but $K$ is
unbounded. Then $K$ contains a ray: there exist $x_0 \in K$ and
$d \neq 0$ with $x_0 + t d \in K$ for all $t \geq 0$. This requires
$\langle a_i, d \rangle \leq 0$ for all~$i$ (otherwise
$\langle a_i, x_0 + td \rangle \to +\infty$, violating
$\langle a_i, x_0 + td \rangle \leq 1$). But $\langle a_i, d \rangle \leq 0$
for all~$i$ contradicts positive spanning. \contradiction

($\Rightarrow$) Suppose positive spanning fails: there exists
$d \neq 0$ with $\langle a_i, d \rangle \leq 0$ for all~$i$.
Since $0 \in \operatorname{int}(K)$, we have $0 + td \in K$ for all
$t \geq 0$: indeed $\langle a_i, td \rangle = t\langle a_i, d\rangle \leq 0 < 1$.
So $K$ contains the ray through~$d$ and is unbounded.
\end{proof}
\end{unverified}


\begin{unverified}
\begin{lemma}[Algorithmic bounded check via triples]\label{lem:bounded-triples}
% [TODO: JÖRN - verify statement and proof]
% Source: check_bounded_f64_first in vertex_enumeration/boundedness.rs
Under the hypotheses of Lemma~\ref{lem:positive-span}, assume
$\operatorname{rank}\{a_1, \ldots, a_F\} = 4$.
Then $\{a_i\}$ positively spans $\mathbb{R}^4$ if and only if for every
triple $\{i,j,k\} \subset \{1,\ldots,F\}$ and every
$d \in \ker(a_i, a_j, a_k) \setminus \{0\}$, there exist $l, m$ with
$\langle a_l, d \rangle > 0$ and $\langle a_m, d \rangle < 0$.
\end{lemma}

\begin{proof}
($\Leftarrow$) If positive spanning fails, there exists $d \neq 0$ with
$\langle a_i, d \rangle \leq 0$ for all~$i$. The dual cone
$C = \{d : \langle a_i, d \rangle \leq 0 \;\forall i\}$ is nontrivial.
Since $\operatorname{rank}\{a_i\} = 4$, the cone $C$ is pointed: if $C$
contained a line (both $d$ and $-d$), then $\langle a_i, d \rangle = 0$
for all~$i$, so $d \in \ker(a_1, \ldots, a_F)^T = \{0\}$, contradicting
$d \neq 0$.

Every pointed polyhedral cone is generated by its extreme rays. An extreme
ray of $C \subset \mathbb{R}^4$ lies on the boundary of at least 3 linearly
independent active constraints $\langle a_i, d \rangle = 0$, so it is
a nonzero element of $\ker(a_i, a_j, a_k)$ for some triple.
For this extreme ray~$d^*$, all remaining dual vertices satisfy
$\langle a_l, d^* \rangle \leq 0$, so the triple check fails.

($\Rightarrow$) If the triple check fails for some $(i,j,k)$ with
kernel direction~$d$, then $\langle a_l, d \rangle \leq 0$ for all
$l \neq i,j,k$ (and $\langle a_i, d \rangle = \langle a_j, d \rangle
= \langle a_k, d \rangle = 0$), so positive spanning fails.
\end{proof}
\end{unverified}


\begin{unverified}
\begin{lemma}[Exact vertex enumeration]\label{lem:vertex-enumeration}
% [TODO: JÖRN - verify statement and proof]
% Source: enumerate_vertices_int in vertex_enumeration/enumerate.rs
Let $K = \{x \in \mathbb{R}^4 : \langle a_i, x \rangle \leq 1 \;\forall i\}$
be a bounded polytope with $0 \in \operatorname{int}(K)$.
A point $v \in \mathbb{R}^4$ is a vertex of~$K$ if and only if there exists
a subset $S \subset \{1,\ldots,F\}$ with $|S| \geq 4$ such that:
\begin{enumerate}
  \item the matrix $(a_s)_{s \in S'}$ is nonsingular for some
    $S' \subset S$ with $|S'| = 4$,
  \item $\langle a_s, v \rangle = 1$ for all $s \in S$ (the defining
    constraints are tight), and
  \item $\langle a_i, v \rangle \leq 1$ for all $i \notin S$ (the
    remaining constraints are satisfied).
\end{enumerate}
The algorithm iterates over all $\binom{F}{4}$ four-element subsets $S'$,
solves $\langle a_{s_k}, v \rangle = 1$ ($k = 1,\ldots,4$) via Cramer's
rule, checks feasibility, and records the full incidence set
$S = \{i : \langle a_i, v \rangle = 1\}$.
\end{lemma}

\begin{proof}
($\Rightarrow$) Let $v$ be a vertex of~$K$. Then $v$ is an extreme point:
$v$ cannot be written as a convex combination of two distinct points in~$K$.
Let $S = \{i : \langle a_i, v \rangle = 1\}$ be the active set.
The active constraints must have $\operatorname{rank}\{a_s : s \in S\} = 4$,
since otherwise $v$ would not be an extreme point (one could perturb $v$
within the nullspace while remaining feasible). So $|S| \geq 4$ and some
4-element subset $S' \subset S$ has nonsingular coefficient matrix.
Feasibility of $v$ gives condition (3).

($\Leftarrow$) If conditions (1)--(3) hold, then $v$ is the unique solution
of the $4 \times 4$ nonsingular system $\langle a_{s_k}, v \rangle = 1$,
so $v$ cannot be written as a convex combination of two distinct feasible
points (they would also need to satisfy the 4 tight constraints, forcing
them to equal~$v$). Hence $v$ is an extreme point, i.e., a vertex.

The algorithm finds all vertices because every vertex has at least one
4-element subset of active constraints with nonsingular coefficient matrix
(by the rank argument above), and the algorithm tests all such subsets.
Deduplication handles non-simple vertices found via multiple subsets.
\end{proof}
\end{unverified}


\begin{unverified}
\begin{lemma}[Irredundancy check]\label{lem:irredundancy}
% [TODO: JÖRN - verify statement and proof]
% Source: check_irredundancy_f64_first in vertex_enumeration/irredundancy.rs
Let $K$ be as above with vertices $V$ and vertex-facet incidence
$V_i = \{v \in V : \langle a_i, v \rangle = 1\}$.
The $i$-th constraint is irredundant if and only if
$\operatorname{aff}(V_i)$ is a hyperplane, equivalently
$\operatorname{affrank}(V_i) = 3$.
\end{lemma}

\begin{proof}
The $i$-th constraint defines the facet
$F_i = K \cap \{x : \langle a_i, x \rangle = 1\}$.
The vertices of $F_i$ are exactly $V_i$ (the vertices of~$K$ that lie
on the hyperplane $\langle a_i, x \rangle = 1$).

If $\operatorname{affrank}(V_i) = 3$, then $F_i = \operatorname{conv}(V_i)$
is a 3-dimensional polytope, so removing the constraint
$\langle a_i, x \rangle \leq 1$ strictly enlarges~$K$ (the interior
of $F_i$ relative to the hyperplane becomes exposed). Hence the
constraint is irredundant.

If $\operatorname{affrank}(V_i) < 3$, then $V_i$ lies in a
2-dimensional affine subspace, so $F_i$ is at most 2-dimensional.
A face of dimension $< \dim(K) - 1 = 3$ is not a facet, meaning
the constraint is redundant: it is implied by the other constraints
in a neighborhood of $F_i$.
\end{proof}
\end{unverified}


\begin{unverified}
\begin{proposition}[Integer-scaled Cramer's rule]\label{prop:integer-cramer}
% [TODO: JÖRN - verify statement and proof]
% Source: optimized vertex enumeration in vertex_enumeration/enumerate.rs
Let $a_1, \ldots, a_F \in \mathbb{Q}^4$ be the dual vertices defining
$K = \{x \in \mathbb{R}^4 : \langle a_i, x \rangle \leq 1\}$.
Let $D = \operatorname{lcm}$ of all denominators of all components of
$a_1, \ldots, a_F$, and set $A_i = D \cdot a_i \in \mathbb{Z}^4$.

For a four-element subset $S = \{s_0, s_1, s_2, s_3\}$, let $M_S$ be the
$4 \times 4$ matrix with rows $A_{s_k}$, and let $\delta = \det(M_S)$.
If $\delta \neq 0$, the candidate vertex~$v$ solving
$\langle a_{s_k}, v \rangle = 1$ for all $k$ has coordinates
\[
  v^{(j)} = \frac{D \cdot \nu_j}{\delta},
  \qquad
  \nu_j = \det\!\bigl(M_S \text{ with column } j
    \text{ replaced by } \mathbf{1}\bigr),
\]
where $\mathbf{1} = (1,1,1,1)^\top$.

The feasibility gap for facet $i \notin S$ is
\[
  1 - \langle a_i, v \rangle
  = \frac{g_i}{\delta},
  \qquad
  g_i = \delta - \sum_{j=0}^{3} A_i^{(j)} \nu_j \;\in\; \mathbb{Z}.
\]
Hence $1 - \langle a_i, v \rangle \geq 0$ if and only if
$\operatorname{sign}(g_i) = \operatorname{sign}(\delta)$ or $g_i = 0$,
and all intermediate computations are in~$\mathbb{Z}$.
\end{proposition}

\begin{proof}
The system $\langle a_{s_k}, v \rangle = 1$ for $k = 0, \ldots, 3$
is equivalent to $(A_{s_k}/D)^\top v = 1$, i.e.,
$M_S\, v = D \cdot \mathbf{1}$.
By Cramer's rule,
\[
  v^{(j)} = \frac{\det(M_S^{(j)})}{\delta},
\]
where $M_S^{(j)}$ is $M_S$ with column~$j$ replaced by
$D \cdot \mathbf{1}$.
By multilinearity of the determinant in the $j$-th column,
$\det(M_S^{(j)}) = D \cdot \nu_j$.

For the feasibility gap:
\[
  1 - \langle a_i, v \rangle
  = 1 - \sum_{j=0}^{3} \frac{A_i^{(j)}}{D} \cdot \frac{D\, \nu_j}{\delta}
  = 1 - \frac{1}{\delta} \sum_{j=0}^{3} A_i^{(j)} \nu_j
  = \frac{\delta - \sum_j A_i^{(j)} \nu_j}{\delta}
  = \frac{g_i}{\delta}.
\]
Since $A_i^{(j)}, \nu_j, \delta \in \mathbb{Z}$, we have $g_i \in \mathbb{Z}$.
The sign condition follows: $g_i / \delta \geq 0$ iff
$\operatorname{sign}(g_i) = \operatorname{sign}(\delta)$ or $g_i = 0$.
\end{proof}
\end{unverified}


\begin{unverified}
\begin{lemma}[Rational arithmetic pipeline]\label{lem:rational-pipeline}
% [TODO: JÖRN - verify statement, add proof]
% Source: rational_arithmetic.rs doc comment
All discrete/combinatorial decisions in the geometry pipeline
(vertex adjacency, symplectic sign $\omega_0(a_i, a_j) \gtrless 0$,
halfspace irredundancy) are computed using exact rational arithmetic
to prevent numerical misclassification.
\end{lemma}
\end{unverified}


% ====================================================================
% Shoelace formula
% ====================================================================

\begin{lemma}[Symplectic shoelace formula]\label{lem:shoelace}
% Jörn: text approved (b7a8bc0) — from \begin{lemma} to \end{proof}
Let $z \colon [0, T] \to \mathbb{R}^4$ be a closed curve with
piecewise-constant velocity:
$\dot{z}(t) = w_k$ on $I_k = (\tau_{k-1}, \tau_k)$ for
$k = 1, \ldots, m$, where $0 = \tau_0 < \cdots < \tau_m = T$.
Then
\[
  \int_0^T \langle -J_0\, \dot{z}(t),\, z(t) \rangle \, dt
  = \sum_{1 \leq j < i \leq m}
    |I_j|\, |I_i|\, \omega_0(w_j, w_i).
\]
\end{lemma}

\begin{proof}
Since $z$ is closed, $\sum_{k=1}^m |I_k|\, w_k = 0$.
For $t \in I_i$:
\[
  z(t) = z(0) + \sum_{j=1}^{i-1} |I_j|\, w_j
    + (t - \tau_{i-1})\, w_i.
\]
Substitute into
$\langle -J_0 w_i, z(t) \rangle$ and integrate over $[0, T]$:
\begin{align*}
\int_0^T \!\langle{-}J_0 \dot{z}, z \rangle \, dt
  &= \sum_{i=1}^m \int_{I_i}
    \!\langle -J_0 w_i,\;
      z(0) + \textstyle\sum_{j < i} |I_j| w_j
      + (t{-}\tau_{i-1}) w_i
    \rangle\, dt \\
  &= \underbrace{
      \bigl\langle -J_0 \bigl(\textstyle\sum_i |I_i| w_i\bigr),\,
      z(0) \bigr\rangle
    }_{= \, 0 \;\text{(closure)}}
  + \sum_{i=1}^m \sum_{j=1}^{i-1}
      |I_i|\, |I_j|\, \omega_0(w_j, w_i) \\
  &\quad
  + \sum_{i=1}^m
      \underbrace{\omega_0(w_i, w_i)}_{= \, 0}\;
      \cdot \tfrac{|I_i|^2}{2}.
\end{align*}
The first term vanishes because
$\sum |I_k| w_k = 0$ (closure of~$z$).
The third term vanishes because
$\omega_0$ is antisymmetric.
\end{proof}


% ====================================================================
% HKO counterexample
% ====================================================================

\begin{unverified}
\begin{theorem}[HKO counterexample to Viterbo's conjecture]%
\label{thm:hko-counterexample}
% [TODO: JÖRN - verify statement, add proof/reference]
% Source: known_polytopes.rs doc comment
% Reference: Haim-Kislev and Ostrover 2024
There exists a 10-facet polytope $K \subset \mathbb{R}^4$
(the HKO pentagon counterexample) with systolic ratio
$\operatorname{sys}(K) > 1$, disproving Viterbo's conjecture
in dimension~$4$.
\end{theorem}
\end{unverified}


% ====================================================================
% f64 pre-filter for vertex enumeration
% ====================================================================
% Source: /tmp/prefilter-bound-simple.tex (Jörn + agent session)
% Status: Jörn is finalizing; will be updated when final version arrives.

\providecommand{\one}{\mathbf{1}}
\providecommand{\emach}{\varepsilon_{\mathrm{mach}}}

\subsection{f64 prefilter for $y^\top A^{-1}\one \le 1$}

All norms are $\ell^2$ / spectral.
Let $\emach = 2^{-53}$, $n = 4$, $\one = (1,1,1,1)^\top$
(so $\|\one\| = 2$).

\subsubsection*{Problem}

Given $A \in \mathbb{Q}^{n \times n}$ and $y \in \mathbb{Q}^n$,
decide whether $y^\top A^{-1}\one \le 1$.
Exact rational arithmetic is expensive.
We want an f64 algorithm that returns
\textsc{true}, \textsc{false}, or \textsc{indeterminate},
where the first two are always correct and the third
triggers a fallback to exact arithmetic.

\subsubsection*{Algorithm}

\begin{enumerate}
\item Round $A$ and $y$ to f64, giving $\hat{A}$ and $\hat{y}$.
\item Compute an SVD of $\hat{A}$, giving computed singular values
  $\hat\sigma_1 \ge \cdots \ge \hat\sigma_n$.
  If $\hat\sigma_n = 0$: return \textsc{indeterminate}.
\item Set $\hat\kappa = \hat\sigma_1 / \hat\sigma_n$.
  If $\emach \hat\kappa > 1/4$: return \textsc{indeterminate}.
\item Solve $\hat{A}\hat{v} \approx \one$ using the SVD factors.
\item Compute $\hat{s} = \hat{y}^\top\hat{v}$ and the tolerance
  $\delta = C \cdot \hat\kappa \cdot \emach
  \cdot \|\hat{v}\| \cdot \|\hat{y}\|$
  with $C = 10^4$.
\item If $\hat{s} > 1 + \delta$: return \textsc{false}
  ($y^\top A^{-1}\one > 1$).\\
  If $\hat{s} < 1 - \delta$: return \textsc{true}
  ($y^\top A^{-1}\one \le 1$).\\
  Else: return \textsc{indeterminate}.
\end{enumerate}

If any step produces $\pm\infty$ or NaN: return
\textsc{indeterminate}.
Steps 1--4 cost $O(n^3)$ and depend only on $A$
(reused across all $y$);
steps 5--6 cost $O(n)$ per $y$.

\subsubsection*{Standard facts used}

\begin{unverified}
\begin{fact}[IEEE~754 rounding model]\label{fact:prefilter-round}
Rounding a real number $x$ to f64 gives $\hat{x}$
with $|\hat{x} - x| \le \emach|x|$.
This is the defining property of IEEE~754 round-to-nearest.
\end{fact}
\end{unverified}

\begin{unverified}
\begin{fact}[Entrywise rounding $\to$ spectral norm]\label{fact:prefilter-cast}
If $\hat{A}$ is obtained by rounding each entry of $A$
to f64, then $|\hat{A}_{ij} - A_{ij}| \le \emach|A_{ij}|$
entrywise (by Fact~\ref{fact:prefilter-round}), and
$\|\hat{A} - A\| \le \sqrt{n}\,\emach\|A\|$.

\emph{Proof.}
$\|\hat{A} - A\|^2 \le \|\hat{A} - A\|_F^2
= \sum_{ij}(\hat{A}_{ij} - A_{ij})^2
\le \emach^2 \sum_{ij} A_{ij}^2
= \emach^2\|A\|_F^2 \le \emach^2 n\|A\|^2$,
where the last step uses
$\|M\|_F^2 = \sum_i \sigma_i(M)^2 \le n\sigma_1(M)^2
= n\|M\|^2$. \qed
\end{fact}
\end{unverified}

% Verified from PDF: GVL4 §8.6, Algorithm 8.6.2 preamble (p. 492). No numbered theorem.
% Was "Higham Ch. 5" but Ch. 5 = Polynomials.
\begin{unverified}
\begin{fact}[SVD backward stability;
  {Golub \& Van~Loan~(2013), \S8.6}]\label{fact:prefilter-svd}
The computed SVD of $B \in \mathbb{R}^{n \times n}$ produces
factors $\hat{U}, \hat{\Sigma}, \hat{V}$ satisfying
$B + E_{\mathrm{svd}} = \hat{U}\hat{\Sigma}\hat{V}^\top$
with $\|E_{\mathrm{svd}}\| \le c_{\mathrm{svd}}\emach\|B\|$,
where $c_{\mathrm{svd}} = O(n)$.
\end{fact}
\end{unverified}

% Verified from PDF: Higham Thm. 8.5 (p. 171) for triangular solve; GVL4 §8.6 for SVD.
% Was "Higham Ch. 3, 5" but Ch. 5 = Polynomials.
\begin{unverified}
\begin{fact}[SVD-based solve backward stability;
  {Higham~(2002), Thm.~8.5; Golub \& Van~Loan~(2013), \S8.6}]\label{fact:prefilter-solve}
Solving $Bx = b$ via computed SVD factors in f64 gives $\hat{x}$
satisfying $(B + E)\hat{x} = b + e$ with
$\|E\| \le p\emach\|B\|$ and $\|e\| \le p\emach\|b\|$,
where $p = O(n)$; $p \le 16$ suffices for $n = 4$.
(This combines the SVD backward error with the rounding
in the back-substitution steps
$\hat{U}^\top b$, $\hat{\Sigma}^{-1}(\cdot)$,
$\hat{V}(\cdot)$.)
\end{fact}
\end{unverified}

% Verified from PDF: GVL4 Corollary 8.6.2 (p. 487) for singular values.
% Eigenvalue version is Corollary 8.1.6 (p. 442).
\begin{unverified}
\begin{fact}[Weyl's singular value perturbation theorem;
  {Golub \& Van~Loan~(2013), Cor.~8.6.2}]\label{fact:prefilter-weyl}
For any $B, B' \in \mathbb{R}^{n \times n}$:
$|\sigma_i(B) - \sigma_i(B')| \le \|B - B'\|$.
\end{fact}
\end{unverified}

% Verified from PDF: Banach lemma is not a standalone lemma in Higham; used inline in
% Ch. 7 proofs (e.g. Thm 7.2). Was "Higham Ch. 2" but Ch. 2 = FP Arithmetic.
\begin{unverified}
\begin{fact}[Banach perturbation lemma;
  {Higham~(2002), Ch.~7, \S7.4}]\label{fact:prefilter-banach}
If $B$ is invertible and $\|B^{-1}\|\|E\| < 1$, then
$B + E$ is invertible and
$\|(B+E)^{-1}\| \le \|B^{-1}\|/(1 - \|B^{-1}\|\|E\|)$.

\emph{Proof.}
$(B+E)^{-1} = (I + B^{-1}E)^{-1}B^{-1}
= \bigl(\sum_{k=0}^\infty (-B^{-1}E)^k\bigr)B^{-1}$.
The series converges when $\|B^{-1}E\| < 1$,
and $\|(B+E)^{-1}\| \le \|B^{-1}\|/(1 - \|B^{-1}\|\|E\|)$
by the geometric series bound. \qed
\end{fact}
\end{unverified}

% Verified: Higham (2002), §3.1, Lemma 3.1 (γₙ bound for inner products).
\begin{unverified}
\begin{fact}[Dot product rounding;
  {Higham~(2002), \S3.1, Lem.~3.1}]\label{fact:prefilter-dot}
For $x, y \in \mathbb{F}^n$, the f64-computed inner product
$\hat{s} = \mathrm{fl}(x^\top y)$ satisfies
$|x^\top y - \hat{s}| \le \gamma_n\, |x|^\top|y|$,
where $\gamma_n = n\emach/(1 - n\emach)$
and $|x|$ denotes componentwise absolute values.
\end{fact}
\end{unverified}

\subsubsection*{Error bound}

% [TODO: JÖRN - The proposition must be restated in terms of the
%   COMPUTABLE $\hat\kappa$, not the true $\kappa = \kappa_2(\hat{A})$.
%   As written, the hypothesis $\emach\kappa \le 1/4$ cannot be checked
%   because $\kappa$ is an algebraic real we don't have.
%
%   Fix: restate with condition $\emach\hat\kappa \le t$ for a threshold
%   $t$ depending on $c_{\mathrm{svd}}$. The proof bridges to $\kappa$
%   internally using Weyl: $\kappa \le \hat\kappa/(1 - c_{\mathrm{svd}}
%   \emach\hat\kappa)$.
%
%   Second consequence: the code computes $\delta = C \hat\kappa \emach
%   \|\hat v\| \|\hat y\|$, but the proposition bounds the error with
%   $C \kappa \emach \|\hat v\| \|\hat y\|$. For $\delta$ to be a valid
%   overestimate, we need $\hat\kappa \ge \kappa$ (or a known bound on
%   $\kappa/\hat\kappa$). Weyl only gives a two-sided bound, so this is
%   not guaranteed. Same fix: restate entirely in terms of $\hat\kappa$.
%
%   Blocker: need concrete $c_{\mathrm{svd}}$ for $n = 4$, from
%   Higham (2002) Ch. 5. This is tracked in tasks/numerics.md under
%   "Solver formal writeup".
%
%   Factor-counting issue (from adversarial Opus review):
%   - The tight bound ``$C < 1400$'' undercounts constant factors.
%     The 6 listed $\lesssim$ factors of 2 miss at least two more:
%     (a) $\|e\| = 2p\emach$ contributes a factor of 2 from
%         $\|\one\| = 2$ that is an exact constant, not a
%         $(1 \pm c\emach)$ factor, so it shouldn't be absorbed
%         into $\lesssim$ but must be tracked explicitly.
%     (b) Substituting $\|\hat A^{-1}\| \lesssim \kappa\|\hat v\|$
%         into eq:prefilter-fwd gives $\kappa\|\hat v\| + \kappa\|\hat v\|
%         = 2\kappa\|\hat v\|$, an untracked factor of 2.
%     Correct tight bound: $C \le 21 \cdot 2^8 = 5376$, not $1344$.
%     The safe $C = 10^4$ remains valid.
%
%   Additional minor gaps (from adversarial review):
%   - The rounding model $|\hat x - x| \le \emach|x|$ assumes normal
%     f64 numbers. Subnormals have larger relative error. Unlikely for
%     our polytope matrices but not formally excluded.
%   - The proposition assumes $A$ is invertible (caller's responsibility).
%   - The code computes $\|\hat v\|$ and $\|\hat y\|$ in f64, introducing
%     $O(\emach)$ relative error not mentioned in the proof. Absorbed by
%     the $C = 10^4$ safety factor.]

\begin{unverified}
\begin{proposition}\label{prop:prefilter-bound}
Let $A$ and $\hat{A}$ both be invertible, with
$\kappa = \kappa_2(\hat{A}) = \|\hat{A}\|\|\hat{A}^{-1}\|$
(the true, not computed, condition number of $\hat{A}$).
If $\emach\kappa \le 1/4$, then:
Write $v = A^{-1}\one$ and let $\hat{v}$, $\hat{s}$
be the f64-computed quantities from the algorithm.
Then
\[
  |y^\top v - \hat{s}|
  \;\le\;
  C \cdot \kappa \cdot \emach
  \cdot \|\hat{y}\| \cdot \|\hat{v}\|
\]
with $C = 10^4$ for $n = 4$
(a tighter accounting gives $C < 1400$; see Remark~\ref{rem:prefilter-constants}).
\end{proposition}
\end{unverified}

\begin{unverified}
\begin{corollary}\label{cor:prefilter-soundness}
Under the hypotheses of Proposition~\ref{prop:prefilter-bound},
the algorithm's three-valued output is sound:
\begin{itemize}
\item If the algorithm returns \textsc{false}
  ($\hat{s} > 1 + \delta$), then $y^\top A^{-1}\one > 1$.
\item If the algorithm returns \textsc{true}
  ($\hat{s} < 1 - \delta$), then $y^\top A^{-1}\one < 1$.
\item If the algorithm returns \textsc{indeterminate},
  no conclusion is drawn.
\end{itemize}
In particular, when the algorithm returns \textsc{false}
for some constraint~$y_i$, the point $A^{-1}\one$ lies
outside~$K$, so the subset cannot yield a vertex.
\end{corollary}

\begin{proof}
By Proposition~\ref{prop:prefilter-bound},
$|y^\top v - \hat{s}| \le \delta$.
If $\hat{s} > 1 + \delta$, then
$y^\top v \ge \hat{s} - \delta > 1$.
If $\hat{s} < 1 - \delta$, then
$y^\top v \le \hat{s} + \delta < 1$.
\end{proof}
\end{unverified}

\noindent
\textit{Why this form?}
The total error is $(p\kappa + 1 + n)\emach\|\hat{y}\|\|\hat{v}\|$:
the $p\kappa$ term comes from the solve (the only step that
amplifies errors), while $1$ and $n$ come from casting and
the dot product (basic rounding, no amplification).
Since $\kappa \ge 1$, we can bound
$p\kappa + 1 + n \le (p + 1 + n)\kappa$,
giving a single factor of $\kappa$ with a larger constant $C$.

\medskip
\noindent
\textit{Why $\hat\kappa \approx \kappa$?}
By Fact~\ref{fact:prefilter-svd}, the computed singular values
$\hat\sigma_i$ are the exact singular values of
$\hat{A} + E_{\mathrm{svd}}$ with
$\|E_{\mathrm{svd}}\| \le c\emach\|\hat{A}\|$.
By Fact~\ref{fact:prefilter-weyl},
$|\hat\sigma_i - \sigma_i(\hat{A})| \le c\emach\|\hat{A}\|
= c\emach\sigma_1(\hat{A})$.
The relative error on $\sigma_n$ is
$O(\emach \cdot \sigma_1/\sigma_n) = O(\emach\kappa)$,
so $\hat\kappa = \kappa(1 + O(\emach\kappa))$.
In the regime $\emach\kappa \le 1/4$,
$\hat\kappa$ and $\kappa$ differ by at most a constant
factor, absorbed into $C$.

\begin{unverified}
\begin{proof}[Proof of Proposition~\ref{prop:prefilter-bound}]
Write $\kappa = \kappa_2(\hat{A})$.
Throughout, $\lesssim$ means $\le$ up to factors of the
form $(1 \pm c\emach)^{\pm 1}$ with $c \le 10^3$.
Since $c\emach < 10^{-13}$, each such factor is $\le 2$.
We number each use; Remark~\ref{rem:prefilter-constants}
enumerates all six.

\medskip\noindent\textbf{Step 1: Backward error of the solve.}
By Fact~\ref{fact:prefilter-solve}, there exist
$E \in \mathbb{R}^{n \times n}$ and
$e \in \mathbb{R}^n$ with
\[
  (\hat{A} + E)\hat{v} = \one + e, \qquad
  \|E\| \le p\emach\|\hat{A}\|, \quad
  \|e\| \le p\emach\|\one\| = 2p\emach,
\]
where $p \le 16$ for $n = 4$.
By Fact~\ref{fact:prefilter-cast},
$\|\hat{A} - A\| \le \sqrt{n}\,\emach\|A\|$.
Since $\sqrt{n} = 2 \le p$, the casting error
is dominated by the solve's backward error.
Writing $\Delta A = (\hat{A} - A) + E$:
\[
  (A + \Delta A)\hat{v} = \one + e, \qquad
  \|\Delta A\| \le (p + \sqrt{n})\emach\|\hat{A}\|
  \lesssim p\emach\|A\|.
  \quad\text{[$\lesssim$ \#2: $\|\hat{A}\| \lesssim \|A\|$]}
\]

\medskip\noindent\textbf{Step 2: Forward error.}
From $Av = \one$:
$v - \hat{v} = A^{-1}(\Delta A\,\hat{v} - e)$.
By Fact~\ref{fact:prefilter-banach}
applied to $A = \hat{A} - (\hat{A} - A)$
with $\|\hat{A}^{-1}\|\|\hat{A} - A\| \le
\sqrt{n}\,\emach\|\hat{A}^{-1}\|\|A\|
\lesssim \sqrt{n}\,\emach\kappa \le 1/2$
(the $\lesssim$ uses $\|A\| \lesssim \|\hat{A}\|$,
[$\lesssim$ \#1]):
$\|A^{-1}\| \le \|\hat{A}^{-1}\|/(1 - 1/2)
= 2\|\hat{A}^{-1}\|$
\quad[$\lesssim$ \#3].
Hence
\begin{equation}\label{eq:prefilter-fwd}
  \|v - \hat{v}\| \le \|A^{-1}\|(\|\Delta A\|\|\hat{v}\| + \|e\|)
  \lesssim
  p\emach\bigl(\kappa\|\hat{v}\| + \|\hat{A}^{-1}\|\bigr).
\end{equation}
From $(\hat{A} + E)\hat{v} = \one + e$
and the reverse triangle inequality:
$\|\hat{v}\| \ge
(\|\one\| - \|e\|)/(\|\hat{A}\| + \|E\|)
= (2 - 2p\emach)/\bigl((1 + p\emach)\|\hat{A}\|\bigr)
\gtrsim 1/\|\hat{A}\|$,
so $1/\|\hat{A}\| \lesssim \|\hat{v}\|$
\quad[$\lesssim$ \#4]
and $\|\hat{A}^{-1}\| = \kappa/\|\hat{A}\|
\lesssim \kappa\|\hat{v}\|$.
Substituting:
\begin{equation}\label{eq:prefilter-fwd2}
  \|v - \hat{v}\| \lesssim p\kappa\emach\|\hat{v}\|.
\end{equation}

\medskip\noindent\textbf{Step 3: Scalar error.}
Decompose:
\[
  y^\top v - \hat{s}
  = \underbrace{y^\top(v - \hat{v})}_{T_1}
  + \underbrace{(y - \hat{y})^\top\hat{v}}_{T_2}
  + \underbrace{(\hat{y}^\top\hat{v} - \hat{s})}_{T_3}.
\]
\emph{$T_1$} (solve error):
$|T_1| \le \|y\|\|v - \hat{v}\|$
by Cauchy--Schwarz.
By Fact~\ref{fact:prefilter-cast} applied to the vector $y$
(the $\sqrt{n}$ factor drops since $\|\cdot\|_F = \|\cdot\|_2$
for vectors),
$\|y - \hat{y}\| \le \emach\|y\|$,
so $\|y\| \le \|\hat{y}\|/(1 - \emach) \lesssim \|\hat{y}\|$
\quad[$\lesssim$ \#5].
By~\eqref{eq:prefilter-fwd2}:
$|T_1| \lesssim p\kappa\emach\|\hat{y}\|\|\hat{v}\|$.

\emph{$T_2$} (casting $y$):
$|T_2| \le \|y - \hat{y}\|\|\hat{v}\|
\le \emach\|y\|\|\hat{v}\|
\lesssim \emach\|\hat{y}\|\|\hat{v}\|$.

\emph{$T_3$} (dot product rounding):
by Fact~\ref{fact:prefilter-dot},
$|T_3| \le \gamma_n\, |\hat{y}|^\top|\hat{v}|$
where $\gamma_n = n\emach/(1 - n\emach) \lesssim n\emach$
\quad[$\lesssim$ \#6].
By Cauchy--Schwarz, $|\hat{y}|^\top|\hat{v}| \le \|\hat{y}\|\|\hat{v}\|$,
so $|T_3| \lesssim n\emach\|\hat{y}\|\|\hat{v}\|$.

\medskip\noindent\textbf{Step 4: Combine.}
Since $\kappa \ge 1$, the $T_2$ and $T_3$ terms
are absorbed:
\[
  |y^\top v - \hat{s}|
  \lesssim (p + 1 + n)\kappa\emach\|\hat{y}\|\|\hat{v}\|.
\]
For $n = 4$, $p \le 16$: pre-factor $= 21$.
With 6 factors of 2: $C \le 21 \cdot 2^6 = 1344 < 10^4$.
\end{proof}
\end{unverified}

\begin{unverified}
\begin{remark}\label{rem:prefilter-constants}
The 6 factors of~2:
\#1~$\|A\| \lesssim \|\hat{A}\|$,
\#2~$\|\hat{A}\| \lesssim \|A\|$,
\#3~Banach lemma,
\#4~lower bound on $\|\hat{v}\|$,
\#5~$\|y\| \lesssim \|\hat{y}\|$,
\#6~$\gamma_n$.
Pre-factor $21$, total $C \le 1344$.

For $\kappa = 1$ and $O(1)$ data:
indeterminate zone $\approx 10^4 \cdot 10^{-16} = 10^{-12}$.
For $\kappa = 10^8$: zone $\approx 10^{-4}$.
\end{remark}
\end{unverified}

% References:
% - N. J. Higham, Accuracy and Stability of Numerical Algorithms,
%   2nd ed., SIAM, 2002.
% - G. H. Golub and C. F. Van Loan, Matrix Computations,
%   4th ed., Johns Hopkins University Press, 2013.
